\section{Failure (Motivation): Accurate Systems With No Visual Headroom}
\label{sec:obs}

On Med-CMR~\cite{gong2025medcmr}, the benchmark of the MedReason~2026 challenge~\cite{medreason2026}, the correct option is the shortest of the five in 91.9\% of 13{,}687 questions (95\% CI 91.4--92.4), so a rule that picks it scores 91.1\% on the test slice.
Three benchmark-tuned model families (the released adapter of a top-3 system; Qwen2.5-VL and InternVL3 models that we fine-tune on the training split and evaluate on 3{,}000 questions they never saw) are \emph{text-saturated}: they stay confident when the image is replaced by a black one for 82--84\% of the questions, and on these \emph{text-decisive} questions no counterfactual image, whether black, random or a plausible wrong one, changes any answer (0 of 2{,}518, 2{,}503 and 2{,}463 questions, one count per system).
Fine-tuning installs the habit (the shortest option is picked 23--25\% of the time before and 91--92\% after), and the headline true-vs-black gaps of the three systems are within 1.3 points, so a conventional ablation calls them equally image-blind.
The stratified view tells them apart: on the 5--15\% text-undecided questions the released adapter gains $+12.1$ points from the image ($p=2\times10^{-5}$; on data that may overlap its unknown training set), whereas our fine-tuned models gain between $-4.1$ and $+2.6$ points across seven training configurations (Supp.~G).
High accuracy is therefore no evidence that these systems use the image, and the average drop without it hides where the image can matter at all.
The finding does not depend on this model-based stratifier: in terciles of a text-only option classifier that uses neither the model nor any image, at most 0.2\% of the answers in the most text-predictable third change under any counterfactual image (Supp.~Z).
Overall accuracy also hides where the cue points away; the released system keeps $83.9\%$ ($[73.2,92.9]$) on the 56 held-out questions whose correct option is not the shortest against $97.1\%$ on the others, our fine-tuned models lose more: on the 206 FRESH questions whose correct option is not the shortest, accuracy is $68.0\%$ (interval $[61.6,74.3]$; Qwen-3B) and $73.3\%$ ($[66.5,79.1]$; Qwen-7B) against $97.5$--$97.7\%$ on the others, and a black image changes it by about a point (Supp.~Z).
A text-only audit of seven multiple-choice sets (Supp.~Z) finds the length cue only in Med-CMR (the shortest-option rule scores 91.9\% against 20\% chance; the other sets are within 6 points of chance), but an option-only classifier that sees neither the question nor the image, the \emph{blind ceiling}, reaches 93.7\% on Med-CMR and 94.1\% on OmniMedVQA against 23--40\% on PMC-VQA, A-OKVQA, MedXpertQA-MM and MMMU health: saturation by text alone is not unique to Med-CMR, but it is far from general; the next section isolates what drives it.
